\documentclass[11pt]{article}

\usepackage[a4paper,margin=1in]{geometry}
\usepackage{amsmath,amssymb,amsthm,mathtools}
\usepackage{mathrsfs}
\usepackage{enumitem}
\usepackage{hyperref}
\usepackage{tikz-cd}
\usepackage{array}

\hypersetup{
	colorlinks=true,
	linkcolor=blue,
	urlcolor=blue,
	citecolor=blue
}

\setlength{\parskip}{0.5em}
\setlength{\parindent}{0pt}

\newtheorem{theorem}{Theorem}[section]
\newtheorem{proposition}[theorem]{Proposition}
\newtheorem{lemma}[theorem]{Lemma}
\newtheorem{corollary}[theorem]{Corollary}
\theoremstyle{definition}
\newtheorem{definition}[theorem]{Definition}
\newtheorem{example}[theorem]{Example}
\newtheorem{exercise}[theorem]{Exercise}
\theoremstyle{remark}
\newtheorem{remark}[theorem]{Remark}

\newcommand{\R}{\mathbb{R}}
\newcommand{\C}{\mathbb{C}}
\newcommand{\Z}{\mathbb{Z}}
\newcommand{\Q}{\mathbb{Q}}
\newcommand{\F}{\mathbb{F}}
\newcommand{\A}{\mathbb{A}}
\newcommand{\PP}{\mathbb{P}}
\newcommand{\cO}{\mathcal{O}}
\newcommand{\cI}{\mathcal{I}}
\newcommand{\cL}{\mathcal{L}}
\newcommand{\Syz}{\mathrm{Syz}}
\newcommand{\im}{\mathrm{im}}
\newcommand{\kerR}{\mathrm{ker}}
\newcommand{\ord}{\mathrm{ord}}
\newcommand{\Spec}{\mathrm{Spec}}
\newcommand{\Hom}{\mathrm{Hom}}
\newcommand{\rank}{\mathrm{rank}}

\title{Lecture Notes on Syzygies\\
	\large A Simple but Penetrating Introduction}
\author{}
\date{}

\begin{document}
	
	\maketitle
	
	\tableofcontents
	
	\section{Why study syzygies?}
	
	In algebra, geometry, and topology, we often begin with \emph{generators}. A module, an ideal, or a ring may be described by a list of elements from which everything else is built. But generators are only the first layer of structure. The next question is:
	
	\begin{center}
		\emph{How are the generators related to one another?}
	\end{center}
	
	A \textbf{syzygy} is exactly such a relation.
	
	This idea is fundamental because mathematics often proceeds in layers:
	\[
	\text{generators} \longrightarrow \text{relations} \longrightarrow \text{relations among relations} \longrightarrow \cdots
	\]
	The word \emph{syzygy} refers to one of these layers of relations.
	
	These notes explain the idea carefully, starting with elementary examples and moving toward free resolutions and geometry.
	
	\section{First definition}
	
	\begin{definition}
		Let \(R\) be a ring and let \(f_1,\dots,f_r \in R\). A \textbf{syzygy} among \(f_1,\dots,f_r\) is a tuple
		\[
		(a_1,\dots,a_r)\in R^r
		\]
		such that
		\[
		a_1f_1+\cdots+a_rf_r=0.
		\]
	\end{definition}
	
	So a syzygy is a relation among the elements \(f_1,\dots,f_r\), where the coefficients \(a_i\) are taken from the ring \(R\).
	
	\begin{remark}
		This looks like linear dependence, but over a ring rather than a field. That is why the theory is richer than ordinary linear algebra.
	\end{remark}
	
	\section{The first key example}
	
	Let
	\[
	R = k[x,y]
	\]
	for a field \(k\). Consider the two elements
	\[
	f_1=x,\qquad f_2=y.
	\]
	A syzygy is a pair \((a,b)\in R^2\) such that
	\[
	ax+by=0.
	\]
	
	\subsection{A visible syzygy}
	
	Take
	\[
	a=-y,\qquad b=x.
	\]
	Then
	\[
	(-y)x + x y = -xy+xy=0.
	\]
	So
	\[
	(-y,x)
	\]
	is a syzygy of \(x\) and \(y\).
	
	This is the first example one should remember.
	
	\subsection{Why it matters}
	
	At first, \(x\) and \(y\) look like two independent generators. But the pair \((-y,x)\) shows that they are tied together by a structured cancellation:
	\[
	(-y)\cdot x + x\cdot y = 0.
	\]
	This is the essence of a syzygy.
	
	\subsection{All syzygies of \(x\) and \(y\)}
	
	Now suppose
	\[
	ax+by=0
	\]
	for some \(a,b\in k[x,y]\). We claim that every such relation is a multiple of \((-y,x)\).
	
	\begin{proposition}
		In \(R=k[x,y]\),
		\[
		\Syz(x,y)=R\cdot(-y,x).
		\]
		That is, every syzygy of \(x\) and \(y\) has the form
		\[
		(-cy,cx)
		\]
		for some \(c\in R\).
	\end{proposition}
	
	\begin{proof}
		Suppose
		\[
		ax+by=0.
		\]
		Then
		\[
		ax=-by.
		\]
		Hence \(ax\) is divisible by \(y\). Since \(x\) and \(y\) are relatively prime in \(k[x,y]\), it follows that \(a\) must be divisible by \(y\). So we may write
		\[
		a=-cy
		\]
		for some \(c\in R\). Substituting this into the relation gives
		\[
		(-cy)x+by=0,
		\]
		so
		\[
		y(-cx+b)=0.
		\]
		Since \(k[x,y]\) is an integral domain, we get
		\[
		-cx+b=0,
		\]
		hence
		\[
		b=cx.
		\]
		Therefore
		\[
		(a,b)=(-cy,cx)=c(-y,x).
		\]
	\end{proof}
	
	\begin{remark}
		This proposition is extremely important conceptually. There may be infinitely many relations, but they are all generated by one basic relation.
	\end{remark}
	
	\section{The module of syzygies}
	
	The previous example suggests that syzygies themselves form an algebraic object.
	
	\begin{definition}
		Let \(R\) be a ring and \(f_1,\dots,f_r\in R\). The set of all syzygies among them is called the \textbf{syzygy module}:
		\[
		\Syz(f_1,\dots,f_r)
		:=
		\left\{
		(a_1,\dots,a_r)\in R^r \;\middle|\; a_1f_1+\cdots+a_rf_r=0
		\right\}.
		\]
	\end{definition}
	
	\begin{proposition}
		\(\Syz(f_1,\dots,f_r)\) is an \(R\)-submodule of \(R^r\).
	\end{proposition}
	
	\begin{proof}
		If
		\[
		(a_1,\dots,a_r),\ (b_1,\dots,b_r)\in \Syz(f_1,\dots,f_r),
		\]
		then
		\[
		\sum a_if_i=0,\qquad \sum b_if_i=0.
		\]
		Hence
		\[
		\sum (a_i+b_i)f_i = \sum a_if_i + \sum b_if_i = 0.
		\]
		So the sum is again a syzygy.
		
		If \(c\in R\), then
		\[
		\sum (ca_i)f_i = c\sum a_if_i = c\cdot 0 = 0.
		\]
		So scalar multiples are again syzygies.
	\end{proof}
	
	\section{Syzygies as kernels}
	
	The cleanest way to understand syzygies is through kernels.
	
	Given \(f_1,\dots,f_r\in R\), define a map
	\[
	\varphi:R^r\to R,\qquad
	(a_1,\dots,a_r)\mapsto a_1f_1+\cdots+a_rf_r.
	\]
	Then
	\[
	\Syz(f_1,\dots,f_r)=\kerR(\varphi).
	\]
	
	Thus syzygies are not an isolated notion: they are simply kernels of natural maps.
	
	\subsection{The example \(x,y\)}
	
	For \(f_1=x,\ f_2=y\), the map is
	\[
	\varphi:R^2\to R,\qquad \varphi(a,b)=ax+by.
	\]
	Then
	\[
	\kerR(\varphi)=\Syz(x,y)=R\cdot(-y,x).
	\]
	
	If we think of the ideal \((x,y)\subseteq R\), then \(\varphi\) lands in \((x,y)\), and we get an exact sequence
	\[
	0\to R\xrightarrow{\psi} R^2\xrightarrow{\varphi} (x,y)\to 0,
	\]
	where
	\[
	\psi(c)=(-cy,cx).
	\]
	
	This exact sequence means:
	
	\begin{itemize}[leftmargin=2em]
		\item \(R^2\) records the two generators \(x,y\),
		\item the map \(R^2\to (x,y)\) sends the standard basis vectors to those generators,
		\item the kernel is exactly the relation among them.
	\end{itemize}
	
	\section{A second example: three generators}
	
	Now let us look at
	\[
	f_1=x^2,\qquad f_2=xy,\qquad f_3=y^2
	\]
	in \(R=k[x,y]\).
	
	We seek triples \((a,b,c)\in R^3\) such that
	\[
	a x^2 + bxy + c y^2 = 0.
	\]
	
	Two obvious syzygies are:
	\[
	(-y,x,0)
	\]
	because
	\[
	(-y)x^2 + x(xy) + 0\cdot y^2 = -x^2y+x^2y=0,
	\]
	and
	\[
	(0,-y,x)
	\]
	because
	\[
	0\cdot x^2 + (-y)(xy) + x y^2 = -xy^2+xy^2=0.
	\]
	
	So the generators \(x^2,xy,y^2\) satisfy at least these two relations.
	
	\begin{remark}
		This example shows that once one has more generators, several basic syzygies may appear.
	\end{remark}
	
	\subsection{A useful perspective}
	
	Notice that \(x^2,xy,y^2\) are not random: they are the degree-\(2\) monomials in two variables. Their syzygies come from the fact that the products
	\[
	x\cdot (xy)=y\cdot (x^2),\qquad
	x\cdot (y^2)=y\cdot (xy)
	\]
	coincide. So syzygies often arise from ``different ways of writing the same product.''
	
	\section{Syzygies versus linear algebra}
	
	It is useful to compare syzygies with ordinary linear dependence.
	
	\subsection{Linear algebra}
	
	If \(v_1,\dots,v_r\) are vectors in a vector space over a field \(k\), a linear relation is
	\[
	c_1v_1+\cdots+c_rv_r=0,\qquad c_i\in k.
	\]
	
	\subsection{Syzygies}
	
	If \(f_1,\dots,f_r\in R\), then a syzygy is
	\[
	a_1f_1+\cdots+a_rf_r=0,\qquad a_i\in R.
	\]
	
	The key difference is that the coefficients \(a_i\) are elements of the ring, not scalars from a field. Thus the coefficients themselves may vary polynomially.
	
	\begin{remark}
		A syzygy is therefore like a ``linear dependence with variable coefficients.''
	\end{remark}
	
	\section{Generators, relations, and presentation}
	
	Suppose \(M\) is an \(R\)-module generated by \(m_1,\dots,m_r\). Then there is a surjective map
	\[
	\pi:R^r\to M,\qquad
	(a_1,\dots,a_r)\mapsto a_1m_1+\cdots+a_rm_r.
	\]
	Its kernel consists of all relations among the generators \(m_1,\dots,m_r\).
	
	\begin{definition}
		The kernel \(\kerR(\pi)\) is called the \textbf{module of first syzygies} of the chosen generators.
	\end{definition}
	
	Thus one has an exact sequence
	\[
	0\to \kerR(\pi)\to R^r\to M\to 0.
	\]
	
	This is called a \textbf{presentation} of \(M\): one describes \(M\) by generators and relations.
	
	\subsection{A simple slogan}
	
	\begin{center}
		\emph{A syzygy is a relation among generators.}
	\end{center}
	
	\section{Higher syzygies}
	
	The first syzygies may themselves have relations.
	
	Suppose
	\[
	0\to K_1\to F_0\to M\to 0
	\]
	with \(F_0\) free. If \(K_1\) is generated by some elements, then we get a surjection
	\[
	F_1\to K_1\to 0
	\]
	from another free module \(F_1\). The kernel of \(F_1\to K_1\) consists of relations among the first syzygies. These are called \textbf{second syzygies}.
	
	Continuing this way, one gets
	\[
	\cdots \to F_2\to F_1\to F_0\to M\to 0.
	\]
	
	\begin{definition}
		A sequence
		\[
		\cdots \to F_2\to F_1\to F_0\to M\to 0
		\]
		with each \(F_i\) free and exact at every stage is called a \textbf{free resolution} of \(M\).
	\end{definition}
	
	Interpretation:
	\[
	\begin{array}{c|c}
		\text{Module} & \text{Meaning} \\
		\hline
		F_0 & \text{generators of }M \\
		F_1 & \text{first syzygies} \\
		F_2 & \text{syzygies among first syzygies} \\
		F_3 & \text{syzygies among second syzygies} \\
		\vdots & \vdots
	\end{array}
	\]
	
	\section{A concrete free resolution}
	
	Let \(R=k[x,y]\) and consider the ideal
	\[
	I=(x,y).
	\]
	Then the map
	\[
	\varphi:R^2\to I,\qquad \varphi(a,b)=ax+by
	\]
	is surjective, and its kernel is generated by \((-y,x)\).
	
	Hence we have an exact sequence
	\[
	0\to R \xrightarrow{\psi} R^2 \xrightarrow{\varphi} I\to 0,
	\]
	where
	\[
	\psi(c)=(-cy,cx).
	\]
	
	This is a free resolution of \(I\).
	
	It says:
	\begin{itemize}[leftmargin=2em]
		\item \(I\) has two generators, namely \(x\) and \(y\),
		\item those two generators satisfy one basic relation,
		\item there are no further nontrivial relations among that relation in this case.
	\end{itemize}
	
	\section{Why syzygies matter in geometry}
	
	Syzygies are not merely algebraic bookkeeping. They carry geometric meaning.
	
	\subsection{Embedded varieties}
	
	Suppose a projective variety \(X\subseteq \PP^n\) is defined by homogeneous polynomials
	\[
	f_1,\dots,f_r.
	\]
	Then the ideal
	\[
	I(X)=(f_1,\dots,f_r)
	\]
	records the equations cutting out \(X\).
	
	But the equations themselves may satisfy relations:
	\[
	a_1f_1+\cdots+a_rf_r=0.
	\]
	These relations are syzygies among the defining equations.
	
	\subsection{Geometric meaning}
	
	The equations tell us which points lie on \(X\).  
	The syzygies tell us how those equations fit together.
	
	Thus:
	\begin{center}
		\emph{Syzygies reveal hidden structure behind the visible equations.}
	\end{center}
	
	\subsection{A guiding picture}
	
	\[
	\text{variety} \longleftrightarrow \text{ideal of equations} \longleftrightarrow \text{syzygies of equations}
	\]
	
	In modern algebraic geometry, syzygies are used to study:
	\begin{itemize}[leftmargin=2em]
		\item embeddings of curves and surfaces,
		\item complexity of defining equations,
		\item canonical curves,
		\item projective normality,
		\item Koszul cohomology.
	\end{itemize}
	
	\section{A simple geometric example}
	
	Consider the parabola in \(\A^2\):
	\[
	y-x^2=0.
	\]
	Its ideal is generated by one equation:
	\[
	I=(y-x^2).
	\]
	Since there is only one generator, there is no nontrivial syzygy among the generators of this ideal.
	
	Now compare with the parametrized curve
	\[
	x=s^2,\qquad y=st,\qquad z=t^2
	\]
	in projective coordinates. Its image in \(\PP^2\) is cut out by
	\[
	y^2-xz=0.
	\]
	Again the ideal is principal, so there is no first syzygy among generators.
	
	But if an ideal is generated by several polynomials, then their interrelations become important. For more complicated curves, these syzygies encode real geometry.
	
	\section{A penetrating interpretation}
	
	Let us return to the example
	\[
	x,\ y \in k[x,y].
	\]
	The syzygy
	\[
	(-y,x)
	\]
	means:
	\[
	(-y)\cdot x + x\cdot y = 0.
	\]
	
	What does this really say?
	
	It says there are \emph{two ways} to build the same product \(xy\):
	\[
	y\cdot x = x\cdot y.
	\]
	The syzygy records the difference between these two expressions, and that difference is zero.
	
	So a syzygy often measures:
	\begin{center}
		\emph{two different constructions that produce the same result.}
	\end{center}
	
	This viewpoint is extremely useful.
	
	\section{Connection with ideals}
	
	If \(I=(f_1,\dots,f_r)\subseteq R\), then the generators \(f_i\) determine a surjection
	\[
	R^r\to I,\qquad e_i\mapsto f_i,
	\]
	where \(e_1,\dots,e_r\) is the standard basis of \(R^r\).
	
	Then the syzygy module is
	\[
	\Syz(f_1,\dots,f_r)=\kerR(R^r\to I).
	\]
	
	Thus the study of ideals naturally leads to the study of syzygies. To know the generators of an ideal is not yet to fully understand the ideal. One also wants to know how the generators interact.
	
	\section{Minimal generators and minimal syzygies}
	
	In practice one often wants the smallest possible generating set.
	
	\begin{definition}
		A generating set \(f_1,\dots,f_r\) of an ideal or module is called \textbf{minimal} if no generator can be omitted.
	\end{definition}
	
	Once one has chosen minimal generators, one may seek minimal generators for the syzygy module. These are called \textbf{minimal first syzygies}. Then one may continue and seek minimal higher syzygies.
	
	This leads to \textbf{minimal free resolutions}, which are among the most refined algebraic invariants of a module.
	
	\section{The word ``syzygy''}
	
	Historically, the word \emph{syzygy} comes from Greek and means something like ``yoked together'' or ``aligned.'' In astronomy it refers to celestial bodies lying in a line.
	
	The mathematical usage is apt: generators in a syzygy are tied together by a relation. They are not free-standing; they are aligned by a constraint.
	
	\section{A short dictionary}
	
	\begin{center}
		\begin{tabular}{>{\raggedright\arraybackslash}p{0.27\textwidth} >{\raggedright\arraybackslash}p{0.63\textwidth}}
			\textbf{Object} & \textbf{Meaning} \\
			\hline
			Generator & A building block from which the module or ideal is produced \\
			Relation & An expression showing that the generators are not independent \\
			Syzygy & A relation among generators \\
			First syzygy & A relation among the original generators \\
			Second syzygy & A relation among the first syzygies \\
			Free resolution & A layered record of generators, relations, relations among relations, etc.
		\end{tabular}
	\end{center}
	
	\section{Exercises}
	
	\begin{exercise}
		Let \(R=k[x,y]\). Show directly that
		\[
		(-y,x)
		\]
		is a syzygy of \(x\) and \(y\).
	\end{exercise}
	
	\begin{exercise}
		Let \(R=k[x,y]\). Prove that every syzygy of \(x\) and \(y\) is a multiple of \((-y,x)\).
	\end{exercise}
	
	\begin{exercise}
		Find two explicit syzygies among
		\[
		x^2,\ xy,\ y^2.
		\]
	\end{exercise}
	
	\begin{exercise}
		Let \(R=k[x,y,z]\). Show that
		\[
		(-y,x,0),\qquad (-z,0,x),\qquad (0,-z,y)
		\]
		are syzygies of \(x,y,z\).
	\end{exercise}
	
	\begin{exercise}
		Let \(M=R/(x,y)\) with \(R=k[x,y]\). Write down a presentation of \(M\) and describe the first syzygy module.
	\end{exercise}
	
	\begin{exercise}
		Suppose \(f_1,\dots,f_r\in R\). Show that \(\Syz(f_1,\dots,f_r)\) is an \(R\)-submodule of \(R^r\).
	\end{exercise}
	
	\section{Solutions sketch}
	
	\subsection*{Exercise 1}
	Compute:
	\[
	(-y)x + x y = -xy + xy = 0.
	\]
	
	\subsection*{Exercise 2}
	Assume \(ax+by=0\). Then \(ax=-by\), so \(y\mid ax\). Since \(x\) and \(y\) are relatively prime, \(y\mid a\). Write \(a=-cy\). Then substitute and deduce \(b=cx\).
	
	\subsection*{Exercise 3}
	Two syzygies are:
	\[
	(-y,x,0),\qquad (0,-y,x).
	\]
	
	\subsection*{Exercise 4}
	Check directly that each given triple annihilates \((x,y,z)\) under the map
	\[
	R^3\to R,\qquad (a,b,c)\mapsto ax+by+cz.
	\]
	
	\subsection*{Exercise 5}
	The quotient \(M=R/(x,y)\) has presentation
	\[
	R^2 \to R \to R/(x,y)\to 0,
	\]
	where \(R^2\to R\) sends \(e_1\mapsto x\), \(e_2\mapsto y\). The kernel is generated by \((-y,x)\).
	
	\section{Final summary}
	
	The idea of syzygy is simple but deep.
	
	\begin{itemize}[leftmargin=2em]
		\item Start with generators.
		\item Ask for relations among them.
		\item Those relations are syzygies.
		\item Then ask for relations among those relations.
		\item This leads to free resolutions.
	\end{itemize}
	
	The basic example
	\[
	(-y)x + x y = 0
	\]
	already contains the whole philosophy: generators may look separate, but in reality they are bound together by hidden structure.
	
	That hidden structure is what syzygies measure.
	
\end{document}